\documentclass[10pt]{article}
\usepackage[T1]{fontenc}
\usepackage[utf8]{inputenc}
\usepackage[french]{babel}
\usepackage{amsmath,amssymb,mathrsfs,mathpazo}
\usepackage{geometry}
\geometry{a4paper,margin=2cm}
\usepackage{graphicx,array,tabularx,tabularray,colortbl,multirow,multicol,booktabs,enumitem,pifont,ulem,fancybox,diagbox,cancel}
\usepackage{pstricks,pstricks-add,pst-plot,pst-tree,pst-node,pst-text,pst-func,pst-math,pst-bspline,pst-3dplot,pst-eucl,pst-solides3d}
\let\pscalculate\undefined
\usepackage{pst-bezier}
\usepackage{tikz,pgfplots,tkz-tab}
\usetikzlibrary{arrows,arrows.meta,patterns,calc,positioning,trees,decorations.markings,shapes.geometric,shapes.misc,shapes.arrows,intersections,decorations.pathmorphing,decorations.pathreplacing,backgrounds,shadows,babel}
\pgfplotsset{compat=1.18}
\usepgfplotslibrary{fillbetween}
\usepackage[np]{numprint}
\usepackage{listings,makecell,pas-tableur,twoopt}
\let\barre\undefined
\usepackage[pstricks]{tablvar}
\definecolor{darkgreen}{RGB}{0,100,0}
\providecommand{\red}{\color{red}}
\providecommand{\blue}{\color{blue}}
\providecommand{\green}{\color{green}}
\providecommand{\black}{\color{black}}
\providecommand{\white}{\color{white}}
\providecommand{\gray}{\color{gray}}
\providecommand{\cyan}{\color{cyan}}
\providecommand{\magenta}{\color{magenta}}
\providecommand{\yellow}{\color{yellow}}
\providecommand{\violet}{\color{violet}}
\providecommand{\orange}{\color{orange}}
\providecommand{\euro}{\texteuro}
\providecommand{\EUR}{\texteuro}
\providecommand{\ieme}{\textsuperscript{e}}
\providecommand{\degres}{\ensuremath{^\circ}}
\providecommand{\degree}{\ensuremath{^\circ}}
\providecommand{\vv}[1]{\overrightarrow{#1}}
\providecommand{\Coord}[1]{#1}
\providecommand{\decosix}{}
\providecommand{\starredbullet}{\ensuremath{\bullet}}
\pagestyle{empty}
\setlength{\parindent}{0pt}
\usepackage{adjustbox}\DeclareUnicodeCharacter{2236}{:}\DeclareUnicodeCharacter{2190}{\ensuremath{\leftarrow}}
\providecommand{\R}{\mathbb{R}}
\providecommand{\N}{\mathbb{N}}
\providecommand{\Z}{\mathbb{Z}}
\providecommand{\Q}{\mathbb{Q}}
\providecommand{\D}{\mathbb{D}}
\providecommand{\C}{\mathbb{C}}
\providecommand{\e}{{\mathrm{e}}}
\providecommand{\vect}[1]{\overrightarrow{#1}}
\providecommand{\vectt}[1]{\overrightarrow{\mathrm{#1}}}
\providecommand{\vvt}[1]{\overrightarrow{\mathrm{#1}}}
\providecommand{\np}[2][]{{#2~#1}}
\providecommand{\vv}[1]{{\overrightarrow{#1}}}
\renewcommand{\d}{{\mathrm{d}}}
\providecommand{\degres}{{^\circ}}
\providecommand{\degre}{{^\circ}}
\providecommand{\blue}{\color{blue}}
\providecommand{\red}{\color{red}}
\providecommand{\enskip}{\quad}
\providecommand{\llbracket}{\lbrack\!\lbrack}
\providecommand{\rrbracket}{\rbrack\!\rbrack}
\providecommand{\euro}{{\text{€}}}
\providecommand{\no}{{\text{n°}}}
\providecommand{\interval}[2]{{[#1;#2]}}
\providecommand{\intervFO}[2]{{[#1;#2[}}
\providecommand{\dotfill}{\ldots}

\providecommand{\ds}{\displaystyle}
\providecommand{\textding}[1]{#1}
\providecommand{\Oij}{\(\mathrm{O};\vec{\imath},\vec{\jmath}\)}
\providecommand{\Oijk}{\(\mathrm{O};\vec{\imath},\vec{\jmath},\vec{k}\)}
\providecommand{\Ouv}{\(\mathrm{O};\vec{u},\vec{v}\)}
\providecommand{\intervFF}[2]{\ensuremath{[#1;#2]}}
\renewcommand{\euro}{\ifmmode\text{\texteuro}\else\texteuro\fi}
\DeclareUnicodeCharacter{00B0}{\ifmmode^{\circ}\else\textdegree\fi}
\DeclareUnicodeCharacter{20AC}{\ifmmode\text{\texteuro}\else\texteuro\fi}
% Énoncés : Éducation nationale. Transcription et corrigés : APMEP, Denis Vergès et contributeurs. Découpage et classement : CoopMaths, Sébastien Lozano, Éric Elter et contributeurs. Adaptation de la présentation : Maths Explorer.
% CC BY-SA : https://coopmaths.fr/www/about/
% https://coopmaths.fr/alea/static/eam/2025/tex/eam_specifique_2025_06_sujet0v1_1.tex
\begin{document}
\noindent\textit{APMEP / CoopMaths. CC BY-SA. Présentation adaptée par Maths Explorer.}\par\medskip
\textbf{Pour cette première partie, aucune justification n'est demandée et une seule réponse est possible par question. Pour chaque question, reportez son numéro sur votre copie et indiquez votre réponse.}

\subsubsection*{Question 1}
L'opération qui permet de calculer 25\% de 480 est :

\textbf{a.} $\dfrac{480}{25 \times 100}$ \hfill \textbf{b.} $25 \times 480 \times 0,1$ \hfill 
\textbf{c.} $\dfrac{480 \times 100}{25}$ \hfill \textbf{d.} $\dfrac{1}{4} \times 480$

\subsubsection*{Question 2}
Voici trois nombres.

$A = \dfrac{1}{5}$ \hfill $B = \dfrac{19}{100}$ \hfill $C = 0,21$

Le classement par ordre croissant de ces trois nombres est :

\textbf{a.} $A < B < C$ \hfill \textbf{b.} $A < C < B$ \hfill 
\textbf{c.} $B < A < C$ \hfill \textbf{d.} $C < B < A$

\subsubsection*{Question 3}
Voici quatre nombres.

$A = \left(\dfrac{1}{5}\right)^2$ \hfill $B = \left(\dfrac{1}{2}\right)^5$ \hfill $C = 0,05$ \hfill $D = \left(\dfrac{1}{3}\right)^3$

Le plus grand de ces quatre nombres est :

\textbf{a.} $A$ \hfill \textbf{b.} $B$ \hfill \textbf{c.} $C$ \hfill \textbf{d.} $D$

\subsubsection*{Question 4}
Un article augmente de 10\% puis il augmente encore de 10\%. Après ces deux augmentations il a augmenté de :

\textbf{a.} $(10\%)^2$ \hfill \textbf{b.} $19\%$ \hfill \textbf{c.} $20\%$ \hfill \textbf{d.} $21\%$

\subsubsection*{Question 5}
Le tiers d'un quart correspond à la fraction :

\textbf{a.} $\dfrac{1}{7}$ \hfill \textbf{b.} $\dfrac{3}{4}$ \hfill \textbf{c.} $\dfrac{1}{3} \times 4$ \hfill \textbf{d.} $\dfrac{1}{12}$

\subsubsection*{Question 6}
On considère $A = 10 + 0,1 + \dfrac{1}{1000}$. On a :

\textbf{a.} $A = \dfrac{20^{- 1}}{1000}$ \hfill \textbf{b.} $A = \dfrac{1}{1000}$ \hfill 
\textbf{c.} $A = 10,101$ \hfill \textbf{d.} $A = 10,110$

\subsubsection*{Question 7}
On considère $A = 10^{10} + 10^{-10}$. $A$ est environ égal à :

\textbf{a.} $10^0$ \hfill \textbf{b.} $0$ \hfill \textbf{c.} $10^{10}$ \hfill \textbf{d.} $100^0$

\subsubsection*{Question 8}
Une durée de 100 minutes correspond à :

\textbf{a.} 1 heure \hfill \textbf{b.} 1,40 heure \hfill 
\textbf{c.} $\dfrac{5}{3}$ heure \hfill \textbf{d.} 2 heures

\bigskip

\begin{minipage}[t]{0.7\textwidth}
\subsubsection*{Question 9}
On considère une droite $D$ représentée ci-contre.

La seule équation pouvant correspondre à l'équation réduite de la droite $D$ est :

\textbf{a.} $y = x + 3$ \hfill \textbf{b.} $y = x - 3$ \hfill 
\textbf{c.} $y = -x + 3$ \hfill \textbf{d.} $y = -x - 3$
\end{minipage}
\hfill
\begin{minipage}[t]{0.29\textwidth}
\vspace{0pt}
\begin{center}
\begin{tikzpicture}[scale=0.8]
\draw[->] (-1,0) -- (3,0) node[below left] {x};
\draw[->] (0,-1.5) -- (0,2) node[below right] {y};
\draw[very thick, color=red] (-1,-1.5) -- (3,1);
\node[below left] at (0,0) {0};
\node[below right, color=red] at (1,-0.25) {$D$};
\end{tikzpicture}
\end{center}
\end{minipage}

\subsubsection*{Question 10}
On considère la fonction $f$ définie pour tout réel $x$ par $f(x) = 7 - \dfrac{1}{2}(x - 3)^2$.

L'image de 3 par la fonction $f$ est égale à :

\textbf{a.} $7 - \dfrac{1}{2}$ \hfill \textbf{b.} $7 - \dfrac{1}{2}(9 + 9)$ \hfill 
\textbf{c.} $7$ \hfill \textbf{d.} $0$

\subsubsection*{Question 11}
Quand on développe $(x - 3)^2$ on obtient :

\textbf{a.} $x^2 + 9$ \hfill \textbf{b.} $x^2 - 9$ \hfill 
\textbf{c.} $x^2 + 6x - 9$ \hfill \textbf{d.} $x^2 - 6x + 9$

\subsubsection*{Question 12}
Voici deux séries de valeurs.

Série A : 1 ; 2 ; 3 \hfill Série B : 0,5 ; 2 ; 100

Une seule de ces affirmations est exacte :

\textbf{a.} Les deux séries ont la même moyenne et la même médiane. \\
\textbf{b.} Les deux séries ont la même moyenne mais pas la même médiane. \\
\textbf{c.} Les deux séries ont la même médiane mais pas la même moyenne. \\
\textbf{d.} Les deux séries n'ont ni la même moyenne ni la même médiane.

\end{document}
